\subsection{Gibbs envelopes}\label{sec:gibbs}

The soft envelopes of \cref{sec:pcis} (thresholds with side conditions) need a product
proposal. When the statistic has overlapping scopes the proposal is a Gibbs measure, and
its sampler and normaliser can only be certified statistically, as in \cref{sec:local}.
This gives one more statistical rule for PCIS arguments.

\begin{definition}[Gibbs-envelope node]\label{def:gibbs-node}
A node with circuit $C_v$ may carry: a statistic $S=\sum_jg_j(x_{A_j})$ given by scopes and
truth tables; an integer $t'$ and a tilt $\theta\ge0$; either a weight vector $w>0$ and
$\alpha<1$, or the claim that $S$ is supermodular; a $P$-proof of $C_v\vdash_P[S\ge t']$, with
$[S\ge t']$ encoded by a fixed counting circuit (e.g.\ a totalizer); and budgets $L$,
$N_{\rm IS}$, $\delta_v$. The verifier checks the proof, recomputes $\kappa(\theta)$ and checks
$\sum_l\kappa_{il}(\theta)w_l\le\alpha w_i$ for all $i$ (or checks supermodularity and samples by
monotone CFTP with abort, \cref{prop:local-cftp}), and runs steps (T2)--(T5) of TG at
$\theta$ with $U=C_v(X)e^{-\theta(S(X)-t')}$.
\end{definition}

\begin{proposition}[Gibbs envelopes]\label{prop:gibbs}
For every circuit and every argument, a Gibbs-envelope node returns an interval containing
$q_v=D(C_v)$ with probability at least $1-\delta_v$; if its proof is missing or wrong, the
same computation with $U'=\min(1,U)$ still returns a valid lower bound. If $C_v$ has such a
node with a proof of size $\poly(|C_v|)$, a certificate with $\alpha\le1-\gamma$ and
$v_C(\theta)=\E_{D_\theta}[C_ve^{-\theta(S-t')}]\ge1/\poly(n)$, then for $0<\eps\le1$ the interval has
ratio at most $1+\eps$ with probability $1-3\delta_v$ after
$L+O\big((\log(1/\delta_v)+\log\log(1/v_C(\theta)))/(\eps^2v_C(\theta))\big)$ samples (with $N_{\rm IS}$
at least the second term), where $L=O((\theta^2m^2/\eps^2+1)\log(1/\delta_v))$ always suffices,
and $L=O((\theta^2n_Sd^2/(\gamma\eps^2)+\theta m/\eps+1)\log(1/\delta_v))$ when the column sums of
$\kappa(\theta)$ are also at most $1-\gamma$.
\end{proposition}

\begin{proof}
For every $C_v$, $\E_{D_\theta}[C_ve^{-\theta(S-t')}]=e^{\theta t'}q_v/Z(\theta)$, because $D_\theta$ has full
support; the proved implication gives $U\le1$. The proof of \cref{thm:local}(i) used nothing
else, and the certificate is recomputed by the verifier, so no argument can make it
false. Without the implication, $U'\le U$ has mean at most $e^{\theta t'}q_v/Z(\theta)$, so the lower
ends of the intervals for $\E U'$ and $\ln Z$ still give a lower bound. Completeness is the
argument of \cref{thm:local}(ii) with the anticoncentration step replaced by the assumption
on $v_C(\theta)$ and the tilt supplied, so no pilot is run. Without a variance bound,
$\bar V\le LR^2/4$ in \cref{lem:local-ti}, so with $\ell=\ln(10/\delta_v)$,
$u\le\frac23\theta R\ell/L+\theta R\sqrt{\ell/(2L)}$, and
with $R\le m$ and $\theta m/\eps\le(\theta m/\eps)^2+1$ this gives the first value of $L$ (which also
makes $h(S_{\max}-s(0))\le\theta m/L\le\eps/10$). With column sums at most $1-\gamma$,
\cref{lem:local-mix}(c) gives $\mathrm{Var}_{D_{\theta'}}S\le V_{\max}=n_Sd^2/(4\gamma)$ for all
$\theta'\le\theta$, and the bound $u\le2\theta\sqrt{8V_{\max}\ell/L}+6\theta R\ell/L$ (for $L\ge\sqrt{\theta R}$)
from the proof of \cref{thm:local}(ii), which does not use $\theta\le\theta_0$, gives the second; for
$\theta\le\theta_0$ it is at most the $L$ of \cref{thm:local}(ii).
\end{proof}

Compared with \cref{thm:local}, the prover contributes three things. (1) The domination is
proved, so $C_v$ need not be syntactically an LTC: an obfuscated or rewired implementation
of a count, or any circuit that implies one, is handled through the proof system.
(2) The tilt and the weights are supplied and only checked, which removes the pilot run and
can certify regimes where the unweighted bound fails. (3) The threshold $t'$ and the
statistic itself are part of the argument, and by \cref{prop:certgap} they, not the
variance of the weights, set the cost: in its instance the one-line proof $C\vdash g$
replaces $K\ge2^{n/2}$ by $K=1$.

\begin{remark}[block certificates]
Sites may be replaced by blocks of $b$ variables, each resampled exactly from its
conditional law by enumerating $2^b$ assignments. The influence of block $B'$ on block $B$
is the largest total variation distance between the conditional laws of $x_B$ when
$x_{-B}$ changes inside $B'$, computable by enumerating $2^{O(dkb)}$ configurations, and the
proof of \cref{lem:local-mix}(b) goes through with blocks in place of sites. This lets an
argument certify strongly coupled clusters of size $O(\log n)$. (Not implemented.)
\end{remark}

\begin{remark}[lower bounds need no mixing certificate]
For any schedule $0=\theta_0<\dots<\theta_L$ and any transitions $T_l$ that preserve
$D_{\theta_l}$ (heat-bath sweeps, for instance), annealed importance sampling \cite{neal2001} is
unbiased for $p$ whether or not the chains mix: draw $x_0\sim D$, $x_l\sim T_l(x_{l-1},\cdot)$, and
weight $\prod_{l=1}^Le^{(\theta_l-\theta_{l-1})S(x_{l-1})}\cdot C(x_L)e^{-\theta_LS(x_L)}$. By Markov's
inequality the median of $r$ independent estimates exceeds $cp$ with probability at most
$(4/c)^{r/2}$ for $c>4$, so the median divided by $c$ is a lower bound valid on every circuit,
without a certificate. As in \cref{sec:lower}, lower bounds need only a proposal; upper
bounds need the certificate, or exact samples.
\end{remark}
